\section{Introduction}

Financial forecasting is inherently uncertain because stock markets are influenced by macroeconomic events, investor sentiment, geopolitical developments, and other sources of noise. In practice, understanding how confident a model is in a prediction is often more valuable than the prediction itself. An incorrect forecast accompanied by high confidence can lead to poor decisions and significant financial losses.

Traditional machine learning models are typically optimised for point prediction accuracy using metrics such as Mean Squared Error (MSE). While these metrics evaluate predictive accuracy, they do not assess whether the model's confidence estimates are reliable. As a result, such models may produce poorly calibrated uncertainty estimates.

The goal of this project is therefore not only to forecast next-day S\&P~500 returns, but also to evaluate the quality of predictive uncertainty. Two approaches are compared:

\begin{enumerate}
    \item A Long Short-Term Memory (LSTM) network with Monte Carlo Dropout for uncertainty estimation.
    \item A Variational Gaussian Process (GP), which provides uncertainty estimates through a principled Bayesian framework.
\end{enumerate}

Both predictive accuracy and uncertainty quality are evaluated using probabilistic metrics such as Negative Log-Likelihood (NLL) and Expected Calibration Error (ECE). The remainder of this paper is structured as follows. Section \ref{sec:background} introduces the theoretical background on Gaussian Processes, variational inference, and MC Dropout. Section \ref{sec:dataset} describes the dataset and preprocessing pipeline. Section \ref{sec:models} details the two model architectures, evaluation metrics, and hyperparameter tuning procedure. Section \ref{sec:results} presents the experimental results. Section \ref{sec:discussion} discusses the findings, and Section \ref{sec:conclusion} concludes the paper.
